\section*{Chapter \ref{ch:b3:green-industrial} --- Green Chemistry and Industrial Processes}

\begin{solution}{exo:b3:green-industrial:1}
Diels--Alder: 100\,\% (an addition). Esterification: $88.0/(60.0 + 46.0) = 83$\,\%. Wittig:
$104.0/(106.0 + 276.0) = 27$\,\%, triphenylphosphine oxide carrying away most of the mass.
\end{solution}

\begin{solution}{exo:b3:green-industrial:2}
$\mathrm{PMI} = 120/8.0 = 15$; $E = \mathrm{PMI} - 1 = 14$.
\end{solution}

\begin{solution}{exo:b3:green-industrial:3}
One wash of a standard load at a given temperature and cleaning result. Detergents
are dosed differently: a concentrated powder needs less per wash, so equal masses
do not deliver the same service.
\end{solution}

\begin{solution}{exo:b3:green-industrial:4}
2-Methyltetrahydrofuran or ethyl acetate for dichloromethane; heptane (with ethyl
acetate) for hexane; for DMF, ethyl acetate, 2-methyltetrahydrofuran or
dimethyl carbonate, or a coupling in water with \omterm{def:b3:colloids:surfactant}{surfactants}, depending on the
reagents.
\end{solution}

\begin{solution}{exo:b3:green-industrial:5}
1.000 mol of acid; 0.800 mol of ester (\qty{88.0}{g/mol}): yield 80\,\%. AE 83.0\,\%. SF $= 152/106 = 1.43$.
$\mathrm{RME} = 70.4/152 = 0.463$; $\mathrm{AE} \times \text{yield}/\mathrm{SF} = 0.830 \times 0.80/1.43 = 0.463$.
\end{solution}

\begin{solution}{exo:b3:green-industrial:6}
Chlorohydrin: $44.0/(28.0 + 71.0 + 74.1) = 25$\,\%; direct oxidation: 100\,\%. One mole of $\ce{CaCl2}$ per mole of
ethylene oxide: $\qty{1.0e6}{g}/\qty{44.0}{g/mol} \times \qty{111.1}{g/mol} = \qty{2.5}{t}$ per tonne.
\end{solution}

\begin{solution}{exo:b3:green-industrial:7}
$3/8 \times \qty{5.88e4}{mol} \times \qty{44.0}{g/mol} = \qty{0.97}{t}$ of $\ce{CO2}$ per tonne of ammonia. For hydrogen, $\ce{CH4 + 2H2O -> CO2 + 4H2}$:
$1/4$ mole per mole, $\qty{5.0e5}{mol}/4 \times \qty{44.0}{g/mol} = \qty{5.5}{t}$ of $\ce{CO2}$ per tonne of hydrogen.
\end{solution}

\begin{solution}{exo:b3:green-industrial:8}
% ledger: dfg:H2O_l, dfh:H2O-l
One kilogram is \qty{500}{mol}. From $\Delta_rG^\circ = \qty{237.1}{kJ/mol}$: $\qty{1.19e5}{kJ} = \qty{32.9}{kWh}$. From $\Delta_rH^\circ = \qty{285.8}{kJ/mol}$: \qty{39.7}{kWh}
(the difference is heat that the surroundings can supply).
\end{solution}

\begin{solution}{exo:b3:green-industrial:9}
$\qty{1.0e6}{g}/\qty{146.0}{g/mol} = \qty{6.85e3}{mol}$ of $\ce{N2O}$, $\qty{0.30}{t}$; $\times 273 = \qty{82}{t}$ of $\ce{CO2}$ equivalent per tonne of adipic acid.
\end{solution}

\begin{solution}{exo:b3:green-industrial:10}
The RME of a step counts the excess reagent as consumed; if the excess is recovered
and fed back, only the part actually lost is consumed, and the overall mass
efficiency is higher than the product of the step values. In the \omterm{def:b3:green-industrial:e-factor}{E-factor} the
recycled stream is not waste, provided the recycling is inside the boundary;
its energy and losses must then be counted.
\end{solution}

\begin{solution}{exo:b3:green-industrial:11}
Before: $20 \times 0.80 = \qty{16}{kg}$ of solvent, 10\,\% lost: \qty{1.6}{kg} of waste. After: \qty{4.0}{kg}, \qty{0.40}{kg} lost. The
complete \omterm{def:b3:green-industrial:e-factor}{E-factor} falls by 1.2 kg per kilogram of product (and the distillation
energy by three quarters).
\end{solution}

\begin{solution}{exo:b3:green-industrial:12}
Per \omterm{def:b3:green-industrial:lca}{functional unit} and with the same boundary, it would report each impact
category separately (climate, land use, water use\dots) rather than one score. A
decision needs the weights given to these categories, the local scarcity of water
and land, the uncertainty of the data, and whether the bio-based feedstock
competes with food.
\end{solution}

\begin{solution}{pb:b3:green-industrial:1}
\textbf{1.} $\ce{C10H14 + C4H6O3 -> C12H16O + C2H4O2}$ (HF); $\ce{C12H16O + H2 -> C12H18O}$ (Raney nickel);
$\ce{C12H18O + CO -> C13H18O2}$ (palladium).

\textbf{2.} $\ce{C10H14 + C4H6O3 + H2 + CO -> C13H18O2 + C2H4O2}$.

\textbf{3.} $134.0 + 102.0 + 2.0 + 28.0 = \qty{266.0}{g/mol}$ of reactants; ibuprofen \qty{206.0}{g/mol}: AE $= 206.0/266.0 = 77$\,\%.

\textbf{4.} Isobutylbenzene, acetic anhydride, ethyl chloroacetate, sodium ethoxide,
aqueous acid ($\ce{H3O+}$), hydroxylamine and water.

\textbf{5.} $134.0 + 102.0 + 122.5 + 68.0 + 19.0 + 33.0 + 2 \times 18.0 = \qty{514.5}{g/mol}$: AE $= 206.0/514.5 = 40$\,\%.

\textbf{6.} Acetic acid, sodium chloride, ethanol, carbon dioxide, water, ammonia (as an
ammonium salt), and aluminium salts from the stoichiometric aluminium chloride.

\textbf{7.} $0.71 + 0.55 + 0.010 + 0.14 + 0.05 = \qty{1.46}{t}$.

\textbf{8.} PMI 1.46; $E = 0.46$.

\textbf{9.} $0.46 - 0.31 = 0.15$.

\textbf{10.} PMI $= 2.0 + 1.5 = 3.5$; $E = 2.5$.

\textbf{11.} Yields below 100\,\%, excess reagents, solvent and catalyst losses, and
work-up materials all add waste that the \omterm{def:b3:green-industrial:atom-economy}{atom economy} ignores.

\textbf{12.} Prevent waste, maximise \omterm{def:b3:green-industrial:atom-economy}{atom economy}, use catalysts rather than
stoichiometric reagents, avoid unnecessary derivatives and steps, save energy.

\textbf{13.} The stoichiometric aluminium chloride, hydrolysed to aluminium waste in the
work-up; hydrogen fluoride is both catalyst and solvent and is recovered and reused.

\textbf{14.} Raney nickel (a heterogeneous hydrogenation catalyst).

\textbf{15.} A palladium phosphine complex: the \omterm{def:b3:homogeneous-catalysis:carbonylation}{carbonylation} chemistry of
\cref{ch:b3:homogeneous-catalysis}.

\textbf{16.} Acetic anhydride transfers one acetyl group; the other half of the
molecule leaves as acetic acid.

\textbf{17.} Hydrogen fluoride is very toxic and corrosive, carbon monoxide toxic and
flammable: closed equipment of resistant materials, detectors and trained
operators.

\textbf{18.} Each step adds heating, cooling, separations and solvent recovery; three
steps need fewer of all of them.

\textbf{19.} $2.5 \times 3000 = \qty{7500}{t}$ a year.

\textbf{20.} $0.15 \times 3000 = \qty{450}{t}$ a year.

\textbf{21.} About \qty{7000}{t} of waste avoided each year.

\textbf{22.} No: the nature of the waste matters (a tonne of salt water is not a tonne of
a persistent toxic compound), and so do energy, emissions and the hazards of the
reagents.

\textbf{23.} The impacts of making the reagents and energy upstream, of emissions to
air and water, and of the product's use and disposal, per \omterm{def:b3:green-industrial:lca}{functional unit}, in
several impact categories.

\textbf{24.} The three-step catalytic route has an \omterm{def:b3:green-industrial:atom-economy}{atom economy} of about 77\,\% (40\,\% for the
six-step route).
\end{solution}
